% titel: Blatt 4, A6: Elektronenkonfiguration und Quantenzahlen
% Teil II: der durchgerechnete Weg.

\subsection*{Rechnung: Konfigurationen}
\begin{align*}
  \mathrm{He}\ (Z=2) &: 1s^2,
    &\quad & \text{Valenz}:\ n=1,\ l=0,\ m_l = 0.\\
  \mathrm{B}\ (Z=5) &: 1s^2\,2s^2\,2p^1,
    &\quad & \text{Valenz}:\ n=2,\ l=1,\ m_l\in\{-1,0,1\}.\\
  \mathrm{S}\ (Z=16) &: 1s^2\,2s^2\,2p^6\,3s^2\,3p^4
    = \mathrm{[Ne]}\,3s^2\,3p^4,
    &\quad & \text{Valenz}:\ n=3,\ l=1,\ m_l\in\{-1,0,1\}.\\
  \mathrm{As}\ (Z=33) &: \mathrm{[Ar]}\,3d^{10}\,4s^2\,4p^3,
    &\quad & \text{Valenz}:\ n=4,\ l=1,\ m_l\in\{-1,0,1\}.\\
  \mathrm{Sb}\ (Z=51) &: \mathrm{[Kr]}\,4d^{10}\,5s^2\,5p^3,
    &\quad & \text{Valenz}:\ n=5,\ l=1,\ m_l\in\{-1,0,1\}.\\
  \mathrm{Pb}\ (Z=82) &: \mathrm{[Xe]}\,4f^{14}\,5d^{10}\,6s^2\,6p^2,
    &\quad & \text{Valenz}:\ n=6,\ l=1,\ m_l\in\{-1,0,1\}.
\end{align*}

\subsection*{Zusammenhang Perioden/Orbitale/Quantenzahlen}
Die Hauptquantenzahl $n$ der äussersten besetzten Schale ist genau die
Periodennummer: He (1. Periode) $\to n=1$; B (2. Periode) $\to n=2$;
S (3. Periode) $\to n=3$; As (4. Periode) $\to n=4$; Sb (5. Periode)
$\to n=5$; Pb (6. Periode) $\to n=6$. Innerhalb der Schale gibt es je
$l=0,1,2,\dots,n-1$ (s, p, d, f) und zu jedem $l$ die $2l+1$
Orientierungen $m_l = -l,\dots,+l$ (z.\,B. drei $p$-Orbitale).

\begin{ergebnis}
\begin{center}
\begin{tabular}{llll}
\toprule
Element & Konfiguration & $n$ & besetzte $l$ \\
\midrule
He  & $1s^2$                        & 1 & s\\
B   & $1s^2\,2s^2\,2p^1$            & 2 & s, p\\
S   & $[\mathrm{Ne}]\,3s^2\,3p^4$   & 3 & s, p\\
As  & $[\mathrm{Ar}]\,3d^{10}4s^2\,4p^3$ & 4 & s, p, d\\
Sb  & $[\mathrm{Kr}]\,4d^{10}5s^2\,5p^3$ & 5 & s, p, d\\
Pb  & $[\mathrm{Xe}]\,4f^{14}5d^{10}6s^2\,6p^2$ & 6 & s, p, d, f\\
\bottomrule
\end{tabular}
\end{center}
\end{ergebnis}

\begin{kontrolle}
Elektronenzahl je Konfiguration addieren (z.\,B. S: $2{+}2{+}6{+}2{+}4 =
16 = Z$). Periodennummer $= n_{\text{max}}$ der Valenzschale stimmt für alle
sechs Elemente; die $l$-Werte folgen aus der Blockzugehörigkeit (s/p-Block
usw.).
\end{kontrolle}

\begin{merksatz}
Konfiguration nach Aufbauprinzip bis $Z$ füllen (Edelgas-Rumpf abkürzen);
Periodennummer $= $ Hauptquantenzahl $n$ der Valenzschale, Zahl der $m_l$-Werte
$= 2l+1$.
\end{merksatz}
